% ============================================================
% CAPÍTULO 1: INTRODUCCIÓN
% ============================================================
\section{Introducción}
\label{sec:introduccion}

El presente documento constituye la Especificación de Requisitos de Software (ERS) del sistema \textbf{AnxioTest}, una plataforma web diseñada para la evaluación del bienestar emocional. Este desarrollo se enmarca en el \textbf{Proyecto de Vinculación con la Colectividad} de la Universidad de las Fuerzas Armadas ESPE, específicamente en el programa de atención a la \textbf{Misión Social Tumiñahi}, que atiende a adultos mayores y comunidades vulnerables.

La ERS describe de manera estructurada y comprensible los requisitos funcionales y no funcionales del sistema, su contexto operativo, las interfaces, las restricciones técnicas y organizativas, así como los criterios de aceptación que permitirán verificar que el producto cumple con lo esperado.

AnxioTest nace de la necesidad de brindar a los beneficiarios de la misión social una herramienta accesible y confiable que les permita conocer, de forma orientativa, su estado emocional mediante cuestionarios validados científicamente (PHQ-9 y BAI). La plataforma está pensada para ser utilizada por personas con poca familiaridad tecnológica, por lo que se prioriza la simplicidad, la accesibilidad visual y auditiva, y la privacidad de los datos.

\subsection{Propósito}
\label{sec:proposito}

El propósito de este documento es definir de forma clara y precisa los requisitos del sistema AnxioTest, de modo que el equipo de desarrollo, los evaluadores y los patrocinadores tengan una base común para comprender, validar y verificar el comportamiento esperado del producto.

La ERS está dirigida a los siguientes grupos de interés:

\begin{itemize}
  \item \textbf{Equipo de desarrollo:} Para guiar la implementación de la interfaz, la lógica de los cuestionarios y el cálculo de resultados.
  \item \textbf{Equipo de pruebas:} Como base para la definición de casos de prueba y criterios de aceptación.
  \item \textbf{Representantes de la Misión Social Tumiñahi y de la ESPE:} Para validar que el sistema responde a las necesidades reales de la comunidad.
  \item \textbf{Usuarios finales:} Para que conozcan las capacidades y limitaciones de la herramienta.
\end{itemize}

La ERS se ha elaborado siguiendo las directrices del estándar \textbf{IEEE 830}, que garantiza la calidad, trazabilidad y verificabilidad de los requisitos.

\subsection{Ámbito del sistema}
\label{sec:ambito}

AnxioTest es una aplicación web estática que permite a los usuarios realizar autoevaluaciones breves de bienestar emocional mediante los cuestionarios estandarizados \textbf{PHQ-9} (depresión) y \textbf{BAI} (ansiedad). El sistema se ejecuta completamente en el navegador, sin necesidad de conexión a internet permanente ni de almacenamiento en servidores, lo que garantiza la privacidad de los datos.

Las funciones principales que ofrece el sistema son:

\begin{itemize}
  \item Mostrar una página de bienvenida con información clara sobre la herramienta y su propósito.
  \item Permitir seleccionar entre dos evaluaciones: PHQ-9 y BAI.
  \item Ejecutar el cuestionario seleccionado con preguntas secuenciales y opciones de respuesta claras.
  \item Calcular automáticamente el puntaje total y clasificar el resultado según criterios clínicos establecidos.
  \item Presentar resultados interpretativos con mensajes adaptados al nivel de severidad y recomendaciones prácticas.
  \item Ofrecer ayuda contextual mediante modales informativos.
  \item Facilitar la lectura de preguntas mediante síntesis de voz (Text-to-Speech).
  \item Permitir cambiar la apariencia visual mediante temas personalizables (claro, oscuro, alto contraste).
\end{itemize}

\textbf{Exclusiones del sistema:}

\begin{itemize}
  \item No incluye registro ni autenticación de usuarios.
  \item No almacena datos en servidores remotos ni en bases de datos.
  \item No ofrece gestión de usuarios ni administración de contenido.
  \item No se integra con sistemas clínicos ni envía datos a terceros.
  \item No sustituye el diagnóstico ni la evaluación profesional de salud mental.
\end{itemize}

\subsection{Personal involucrado}
\label{sec:personal_involucrado}

El desarrollo de AnxioTest se lleva a cabo bajo la supervisión del \textbf{Ing. Cristian Cola}, docente de la Universidad de las Fuerzas Armadas ESPE, y cuenta con la colaboración de estudiantes de Ingeniería de Software, organizados en tres equipos:

\begin{longtable}{p{0.20\linewidth} p{0.30\linewidth} p{0.45\linewidth}}
\toprule
\rowcolor{tableheader!20}
\textbf{Equipo} & \textbf{Integrantes} & \textbf{Responsabilidades} \\
\midrule
\endfirsthead
\rowcolor{tableheader!20}
\textbf{Equipo} & \textbf{Integrantes} & \textbf{Responsabilidades} \\
\midrule
\endhead
\bottomrule
\endfoot
\textbf{Desarrollo} & Estudiantes de Ingeniería de Software – ESPE & Implementación del frontend (HTML, CSS, JavaScript), lógica de cuestionarios, cálculo de resultados, gestión de almacenamiento local y pruebas unitarias. \\
\textbf{Pruebas} & Estudiantes colaboradores del proyecto & Diseño y ejecución de casos de prueba, pruebas de usabilidad, pruebas de accesibilidad, validación de requisitos y reporte de incidentes. \\
\textbf{Documentación} & Estudiantes colaboradores del proyecto & Elaboración y mantenimiento de la ERS, manuales de usuario, actas de reuniones, y gestión de la trazabilidad de requisitos. \\
\bottomrule
\caption{Equipos de trabajo del proyecto AnxioTest}
\label{tab:equipos}
\end{longtable}

Además, el proyecto cuenta con el apoyo de la \textbf{Misión Social Tumiñahi}, que colabora en la validación de requisitos y en las pruebas de campo con usuarios reales, garantizando que el sistema se ajuste a las necesidades y características de la población objetivo.

\subsection{Definiciones, acrónimos y abreviaturas}
\label{sec:definiciones}

\subsubsection{Definiciones}

\begin{longtable}{p{0.25\linewidth} p{0.70\linewidth}}
\toprule
\rowcolor{tableheader!20}
\textbf{Término} & \textbf{Definición} \\
\midrule
\endfirsthead
\rowcolor{tableheader!20}
\textbf{Término} & \textbf{Definición} \\
\midrule
\endhead
\bottomrule
\endfoot
\textbf{AnxioTest} & Plataforma web de evaluación de bienestar emocional, cuyo nombre combina ``anxiety'' (ansiedad) y ``test'' (prueba). \\
\textbf{Tamizaje} & Proceso de identificación preliminar de posibles trastornos mentales mediante instrumentos estandarizados, con el fin de derivar a evaluación profesional. \\
\textbf{Síntoma} & Manifestación subjetiva de una alteración en el estado de salud reportada por el paciente, indicativa de un posible trastorno. \\
\textbf{Severidad} & Grado de intensidad o gravedad de los síntomas reportados, clasificado en categorías desde mínimo hasta grave. \\
\textbf{Accesibilidad} & Conjunto de características que permiten el uso del sistema por personas con diversas capacidades, incluyendo dificultades visuales, motoras o cognitivas. \\
\textbf{SPA} & Single Page Application. Arquitectura web que carga una única página HTML y actualiza dinámicamente el contenido sin recarga completa. \\
\bottomrule
\caption{Definiciones}
\label{tab:definiciones}
\end{longtable}

\subsubsection{Acrónimos}

\begin{longtable}{p{0.20\linewidth} p{0.75\linewidth}}
\toprule
\rowcolor{tableheader!20}
\textbf{Acrónimo} & \textbf{Descripción} \\
\midrule
\endfirsthead
\rowcolor{tableheader!20}
\textbf{Acrónimo} & \textbf{Descripción} \\
\midrule
\endhead
\bottomrule
\endfoot
\textbf{ERS / SRS} & Especificación de Requisitos de Software. Documento que describe formalmente los requisitos de un sistema software. \\
\textbf{IEEE} & Institute of Electrical and Electronics Engineers. Organización que desarrolla estándares para la ingeniería de software. \\
\textbf{BAI} & Beck Anxiety Inventory. Inventario de Ansiedad de Beck, instrumento de 21 ítems para medir síntomas de ansiedad. \\
\textbf{PHQ-9} & Patient Health Questionnaire -- 9 items. Cuestionario de 9 preguntas para detectar síntomas depresivos. \\
\textbf{TTS} & Text-to-Speech. Tecnología que convierte texto en voz sintetizada. \\
\textbf{UI} & User Interface. Interfaz de usuario. \\
\textbf{UX} & User Experience. Experiencia de usuario. \\
\textbf{ARIA} & Accessible Rich Internet Applications. Atributos HTML que mejoran la accesibilidad. \\
\textbf{WCAG} & Web Content Accessibility Guidelines. Pautas de accesibilidad al contenido web del W3C. \\
\textbf{GDPR} & General Data Protection Regulation. Reglamento General de Protección de Datos. \\
\textbf{RF} & Requisito Funcional. \\
\textbf{RNF} & Requisito No Funcional. \\
\bottomrule
\caption{Acrónimos}
\label{tab:acronimos}
\end{longtable}

\subsubsection{Abreviaturas}

\begin{longtable}{p{0.20\linewidth} p{0.75\linewidth}}
\toprule
\rowcolor{tableheader!20}
\textbf{Abreviatura} & \textbf{Descripción} \\
\midrule
\endfirsthead
\rowcolor{tableheader!20}
\textbf{Abreviatura} & \textbf{Descripción} \\
\midrule
\endhead
\bottomrule
\endfoot
\textbf{ESPE} & Universidad de las Fuerzas Armadas. \\
\textbf{W3C} & World Wide Web Consortium. \\
\textbf{CSS} & Cascading Style Sheets (hojas de estilo en cascada). \\
\textbf{HTML} & HyperText Markup Language. \\
\textbf{JS} & JavaScript. \\
\bottomrule
\caption{Abreviaturas}
\label{tab:abreviaturas}
\end{longtable}

\subsection{Referencias}
\label{sec:referencias}

\begin{longtable}{p{0.06\linewidth} p{0.32\linewidth} p{0.36\linewidth} p{0.12\linewidth}}
\toprule
\rowcolor{tableheader!20}
\textbf{Ref.} & \textbf{Título} & \textbf{Fuente} & \textbf{Fecha} \\
\midrule
\endfirsthead
\rowcolor{tableheader!20}
\textbf{Ref.} & \textbf{Título} & \textbf{Fuente} & \textbf{Fecha} \\
\midrule
\endhead
\bottomrule
\endfoot
{[1]} & IEEE Std 830-1998 --- Recommended Practice for SRS & \url{https://ieeexplore.ieee.org} & 1998 \\
{[2]} & Kroenke et al. The PHQ-9: Validity of a Brief Depression Severity Measure & \textit{J. General Internal Medicine}, 16(9), 606--613 & 2001 \\
{[3]} & Beck et al. An Inventory for Measuring Clinical Anxiety & \textit{J. Consulting and Clinical Psychology}, 56(6), 893--897 & 1988 \\
{[4]} & WCAG 2.1 --- Web Content Accessibility Guidelines & \url{https://www.w3.org/TR/WCAG21/} & 2018 \\
{[5]} & Documentación interna del proyecto AnxioTest, ESPE & Repositorio interno del proyecto & 2025 \\
{[6]} & Lineamientos del Proyecto de Vinculación con la Colectividad, ESPE & Departamento de Vinculación & 2025 \\
{[7]} & Ley Orgánica de Protección de Datos Personales del Ecuador & \url{https://www.registrosdatospublicos.gob.ec} & 2021 \\
\bottomrule
\caption{Referencias}
\label{tab:referencias}
\end{longtable}

\subsection{Visión general del documento}
\label{sec:vision_general}

Este documento de Especificación de Requisitos de Software se organiza siguiendo la estructura recomendada por el estándar IEEE 830:

\begin{itemize}
  \item \textbf{Capítulo 1 -- Introducción:} Presenta el contexto del proyecto, el propósito, el alcance, el personal involucrado, las definiciones y las referencias.
  \item \textbf{Capítulo 2 -- Descripción general del producto:} Ofrece una visión global del sistema: perspectiva, funciones, características de los usuarios, restricciones, suposiciones y evolución previsible.
  \item \textbf{Capítulo 3 -- Requisitos específicos:} Detalla los requisitos funcionales mediante casos de uso (Nivel 0 y Nivel 1), con sus respectivas descripciones y trazabilidad.
  \item \textbf{Capítulo 4 -- Requisitos no funcionales:} Describe los atributos de calidad del sistema (rendimiento, seguridad, usabilidad, accesibilidad, portabilidad, mantenibilidad, disponibilidad) con métricas cuantitativas.
  \item \textbf{Capítulo 5 -- Apéndices:} Contiene información complementaria: glosario, formatos de datos, resúmenes de referencia rápida, declaración de cumplimiento normativo e historial de versiones.
\end{itemize}

Cada requisito (funcional o no funcional) se identifica con un código único (\textbf{RF-XX} o \textbf{RNF-XX}), lo que permite su trazabilidad, verificación y gestión de cambios a lo largo del ciclo de vida del proyecto.